\mspage{2}
\gdef\currentbab{عرض مؤلف}
\pstart\centering بسم الله الرحمن الرحيم\par\pend
\pstart\centering\large\bfseries\urdufont عرضِ مؤلف\par\pend
\begin{urduprose}
\pstart\fontsize{9.4}{18.2}\selectfont
علمائے كرام و بزرگانِ دين كى خدمت ميں يہ عرض ہے كہ ہندوستان كے مدارس ميں ابتدائى عربى صرف و نحو سكھانے كے دو طريقے رائج ہيں، ايک يہ كہ صرف و نحو كو عربى عبارتوں ميں پڑھانا ہے، اس ميں مشكل يہ ہے كہ طلبہ زبانِ عربى سے واقف نہيں تو عربى عبارت كيسے سمجھيں؟ خصوصاً تعليلات طلبہ كے لئے بہت مشكل ہوتى ہيں۔ دوسرا طريقہ فارسى ميں پڑھانا ہے، اس ميں يہ بات ہے كہ طالب علم پہلے اردو سيكھے پھر فارسى پھر عربى سيكھے، پھر بھى كامل طور پر طلبہ اخذ نہيں كر پاتے، خصوصاً بيرونى ممالک كے طلبہ كے لئے يہ طريقہ تقريباً محال ہے۔
\pend
\pstart\fontsize{9.4}{18.2}\selectfont
دوسرى بات يہ ہے كہ عموماً طلبہ كے اذہان ميں يہ بات بيٹھى ہوئى ہے كہ تركيب بہت ہى مشكل ہے، ان سب امور كے پيشِ نظر عربى عبارتوں كے بغير آسان انداز ميں يہ رسالہ پيش كيا گيا ہے كہ ہر ايک اپنى زبان ميں پڑھ لے اور قوانين خود بخود سمجھ ميں آجائيں يا استاذ كے اشارہ سے سمجھ لے اور تركيب بھى سہل انداز ميں پيش كى گئى ہے، پھر آخر ميں جب عربى كى استعداد پيدا ہوجاتى ہے تو عربى عبارتيں بھى دى گئى ہيں، ليكن يہ بات ضرور ہے كہ اساتذہ طلبہ سے جتنى مشق كرائيں گے يہ رسالہ اتنا ہى مفيد ہوگا اور طلبہ صرف و نحو كى ابتدائى باتوں كو خوب سمجھ ليں گے اور آگے بڑھنے كا واضح طريقہ كُھل جائے گا، اس كے بعد طلبہ كو ”ہدايۃ النحو“ يا ”كافيہ“ پڑھا سكتے ہيں۔ آپ حضرات كى خدماتِ عاليہ ميں گذارش ہے كہ اس ميں كوئى غلطى ہو تو براہِ كرم اس كى نشاندہى فرمائيں تاكہ آئندہ اس كى تصحيح كى جاسكے۔
\pend
\pstart
\begin{minipage}[t]{.45\linewidth}والسلام\end{minipage}\hfill
\begin{minipage}[t]{.45\linewidth}\raggedleft\setlength{\baselineskip}{1.8em}
محمد ميران باقوى\\ فاضل ديوبند\\
۸ شوال ۱۳۹۶ھ مطابق ۳ اكتوبر ۱۹۷۶ء\end{minipage}
\pend
\end{urduprose}
